\clearpage\fancyhead[L]{\small JUNO.}\hypertarget{native-330-lower}{}\label{native-330-lower}
\begin{center}\textbf{JUNO.}\par\vspace{4pt}\hrule\vspace{4pt}\fontsize{10}{12}\selectfont Göttingische gelehrte Anzeigen. Stück 92. 93. S. 913\,..\,914. 1811. Juni 10.\par\vspace{4pt}\hrule\vspace{4pt}\end{center}

\par\indent Die letzten Nachrichten über die Juno finden sich im 136. Stück unserer Anz. vom Jahre 1808, [S. 305 d. B.] woselbst die auf der hiesigen Sternwarte von Hrn. Prof. \textsc{Gauss} im Jahre 1808 angestellten Beobachtungen dieses Planeten, die Bestimmung der vierten Opposition, und die achten Elemente mitgetheilt sind. Die fünfte, im Januar 1810 eingetretene, Opposition ist nirgends beobachtet: das erste Mal, dass durch eine allgemeine Versäumniss aller Astronomen die Bestimmung der Opposition eines der vier neuen Planeten verloren gegangen ist. Um so wichtiger war es, zu verhüten, dass auch die \SourceLeaf{331}\emph{sechste}, im April d. J. einfallende,  Opposition unbeobachtet bleibe, und die wenigen, bei dem äusserst geringen Lichte des Planeten nur am Kreismicrometer von Hrn. Prof. \textsc{Gauss} angestellten, Beobachtungen werden demnach desto schätzbarer sein, da sie allem Anschein nach die einzigen sein werden, welche dieses Mal irgendwo gemacht worden sind. Zur Vergleichung wurden $\varphi$ Librae und einige andere in der Nähe befindliche Sterne der Histoire céleste angewandt, nach deren künftiger schärferer Bestimmung die folgenden Resultate noch einer kleinen Berichtigung bedürftig sein werden.
\Needspace{10\baselineskip}
\begin{center}\fontsize{10}{12}\selectfont\setlength{\tabcolsep}{4pt}\renewcommand{\arraystretch}{1.15}\begin{tabular}{c|c|c|c}
1811 & \shortstack{mittlere Zeit\\in Göttingen} & \shortstack{Gerade\\Aufsteigung} & \shortstack{Südliche\\Declination}\\\hline
April 22 & $9^{\text{u}}\;51^{\text{m}}\;35^{\text{s}}$ & $216^\circ\;41'\;50''1$ & $0^\circ\;58'\;16''$\\
24 & $10\;32\;55$ & $216\;17\;58.7$ & $0\;46\;50$\\
25 & $10\;19\;22$ & $216\;6\;3.6$ & $0\;40\;4$\\
26 & $10\;17\;30$ & $215\;54\;7.3$ & $0\;33\;51$\\
\end{tabular}\end{center}

\par\indent Die Declination vom 22. April ist zweifelhaft, und die vom 24. April auch nicht so zuverlässig, als die beiden folgenden, wo der Planet eine bequemere Lage hatte.

\Needspace{8\baselineskip}\par\indent Für die Opposition hat Hr. Prof. \textsc{Gauss} aus diesen Beobachtungen folgendes Resultat gefunden:
\begin{center}1811 April 24. $19^{\text{u}}\;20^{\text{m}}\;12^{\text{s}}$ mittl. Zeit in Göttingen\\$214^\circ\;8'\;48''3$ wahre heliocentrische Länge\\$12\;55\;2.0$ wahre geocentrische Breite, nordl.\end{center}
\Needspace{20\baselineskip}\par\indent In der starken Abweichung dieses Ortes von den auf die vier ersten Oppositionen gegründeten Elementen erkennt man nunmehr auch bei der Juno den Einfluss der Störungen, welche besonders der Jupiter ausübt, und deren Berechnung bei der Juno eine eben so ungeheure Arbeit erfordern wird, wie bei der Pallas: es ist billig, dass diese Arbeit bei der früher entdeckten Pallas zuerst beseitigt sein müsse, daher Hr. Prof. \textsc{Gauss} sich einstweilen begnügt hat, neue elliptische Elemente auf die Oppositionen von 1806, 1807, 1808, 1811, zu gründen, welche wir hier folgen lassen:

\begin{center}\fontsize{10}{12}\selectfont\setlength{\tabcolsep}{4pt}\renewcommand{\arraystretch}{1.15}\begin{tabular}{llr}
\shortstack[l]{Epoche der mittlerm Länge, 1811,\\Meridian von Göttingen} & .\;. & $177^\circ\;48'\;1''8$\\
Tägliche mittlere tropische Bewegung & .\;.\;.\;.\;.\;.\;.\;.\;.\;. & $813''2486$\\
Länge der Sonnennähe 1811 & .\;.\;.\;.\;.\;.\;.\;.\;.\;.\;.\;.\;.\;. & $53^\circ\;14'\;32''4$\\
Länge des aufsteigenden Knoten 1811 & .\;.\;.\;.\;.\;.\;.\;.\;.\;. & $171\;9\;13.5$\\
Neigung der Bahn & .\;.\;.\;.\;.\;.\;.\;.\;.\;.\;.\;.\;.\;.\;.\;.\;. & $13\;4\;27.0$\\
Excentricität & .\;.\;.\;.\;.\;.\;.\;.\;.\;.\;.\;.\;.\;.\;.\;.\;. & $=\sin14\;44\;9.1$\\
Logarithm der halben grossen Axe & .\;.\;.\;.\;.\;.\;.\;.\;.\;.\;. & $0.4265711.$\\
\end{tabular}\end{center}
\clearpage\fancyhead[L]{\small GAUSS AN VON LINDENAU.}\hypertarget{native-331-lower}{}\label{native-331-lower}
\begin{center}\textbf{GAUSS AN VON LINDENAU.}\par\vspace{4pt}\hrule\vspace{4pt}\fontsize{10}{12}\selectfont Monatliche Correspondenz. Band XXIV. Seite 180\,..\,182. 1811. August.\par\vspace{4pt}\hrule\vspace{4pt}\end{center}

\begin{flushright}Göttingen, den 3. August 1811.\end{flushright}
\Needspace{16\baselineskip}\par\indent Ich habe in diesen Tagen aus den Beobachtungen des Cometen, die Herr \textsc{von Zach} angestellt hat, und die Sie an Herrn Prof. \textsc{Harding} geschickt hatten, vorläufig folgende parabolische Elemente abgeleitet:
\SourceLeaf{332}
\begin{center}\fontsize{10}{12}\selectfont Durchgang durch die Sonnennähe 1811 Sept. 10\\$0^{\text{u}}\;51^{\text{m}}\;21^{\text{s}}$ Merid. von Göttingen.\end{center}

\begin{center}\fontsize{10}{12}\selectfont\setlength{\tabcolsep}{4pt}\renewcommand{\arraystretch}{1.15}\begin{tabular}{llr}
Logarithm. des Abstandes in der Sonnennähe & .\;.\;.\;.\;.\;.\;. & $9.99153$\\
Länge der Sonnennähe & .\;.\;.\;.\;.\;.\;.\;.\;.\;.\;.\;.\;.\;.\;. & $73^\circ\;14'\;35''$\\
Länge des aufsteigenden Knoten & .\;.\;.\;.\;.\;.\;.\;.\;.\;.\;.\;. & $141\;4\;59$\\
Neigung der Bahn & .\;.\;.\;.\;.\;.\;.\;.\;.\;.\;.\;.\;.\;.\;.\;.\;. & $73\;48\;2$\\
Bewegung rückläufig &  & \\
\end{tabular}\end{center}

\Needspace{15\baselineskip}\par\indent Nach der Mitte dieses Monats kann man schon anfangen, diesen merkwürdigen Weltkörper am nordlichen Himmel aufzusuchen; ich habe zu dem Ende nach obigen Elementen folgende Positionen berechnet:

\begin{center}\fontsize{10}{12}\selectfont\setlength{\tabcolsep}{4pt}\renewcommand{\arraystretch}{1.15}\begin{tabular}{lrr}
1811 Aug. 11 $12^{\text{u}}$ & AR. $142^\circ\;43'$ & Decl. $29^\circ\;54'$ N.\\
26 & $152\;19$ & $36\;15$\\
Sept. 10 & $167\;34$ & $43\;14$\\
\end{tabular}\end{center}

\begin{center}\fontsize{10}{12}\selectfont\setlength{\tabcolsep}{4pt}\renewcommand{\arraystretch}{1.15}\begin{tabular}{lrr}
1811 Sept. 25 $12^{\text{u}}$ & AR. $194^\circ\;22'$ & Decl. $48^\circ\;35'$ N.\\
Oct. 10 & $232\;28$ & $44\;52$\\
25 & $262\;49$ & $30\;42$\\
\end{tabular}\end{center}

\Needspace{20\baselineskip}\par\indent Diese Resultate können indess nur als vorläufige gelten, da die bisher langsame heliocentrische Bewegung (vom 11. April bis 2. Juni nur $15^\circ$) eine sehr scharfe Bestimmung der Elemente noch nicht erlaubt. Die Lichtstärke (die in der Distanz 1 von Erde und Sonne zur Einheit genommen) war

\begin{center}\fontsize{10}{12}\selectfont\setlength{\tabcolsep}{4pt}\renewcommand{\arraystretch}{1.15}\begin{tabular}{llr}
April 11 & .\;.\;.\;.\;. & $0.035$\\
Mai 7 & .\;.\;.\;.\;. & $0.042$\\
Juni 2 & .\;.\;.\;.\;. & $0.052$\\
\end{tabular}\end{center}

\noindent und wird sein

\begin{center}\fontsize{10}{12}\selectfont\setlength{\tabcolsep}{4pt}\renewcommand{\arraystretch}{1.15}\begin{tabular}{llr}
Aug. 11 & .\;.\;.\;.\;.\;.\;.\;.\;. & $0.194$\\
26 & .\;.\;.\;.\;.\;.\;.\;.\;. & $0.294$\\
Sept. 10 & .\;.\;.\;.\;.\;.\;.\;.\;. & $0.434$\\
Sept. 25 & .\;.\;.\;.\;.\;.\;.\;.\;. & $0.576$\\
Oct. 10 & .\;.\;.\;.\;.\;.\;.\;.\;. & $0.595$\\
25 & .\;.\;.\;.\;.\;.\;.\;.\;. & $0.415$\\
\end{tabular}\end{center}

\noindent Erst gegen den Februar wird er im Wassermann unsichtbar, und vielleicht im April oder Mai daselbst von neuem sichtbar werden, falls sein Licht nicht gar zu schwach ist. Für den 2. Juli 1812 habe ich flüchtig berechnet:
\begin{center}AR. $336^\circ$\quad Decl. $22^\circ$ südl.\quad Lichtstärke $0.005$.\end{center}
\par\indent [Handschriftl. Bemerk. Nach verbesserten Elementen 1812 Juni $3.22^{\mathrm u}$ $337^\circ\,49'$ --- $11^\circ\,41'$ --- $0.5435$ $0.0057$.]

\par\indent Die Störungen der \emph{Pallas} habe ich nun bis Ende 1812 fortgesetzt, und in einigen Tagen hoffe ich Ihnen die Ephemeride, welche jetzt von Hrn. \textsc{Nicolai}, einem sehr geübten jungen Mann berechnet wird, schicken zu können.
